\documentclass[12pt]{article}
\usepackage{lmodern}
\usepackage{geometry}
\usepackage{enumitem}
\setlist[description]{leftmargin=1cm,labelindent=1cm}
\usepackage{amsmath}
\usepackage{amsthm}
\usepackage{amssymb}
\usepackage{mathtools}
\usepackage[portuguese]{babel}

\theoremstyle{plain}
\newtheorem{thm}{Teorema}[section]
\newtheorem{prop}[thm]{Proposição}
\newtheorem{lem}[thm]{Lema}
\newtheorem{cor}[thm]{Corolário}

\theoremstyle{definition}
\newtheorem{defn}[thm]{Definição}
\newtheorem{exmp}[thm]{Exemplo}
\newtheorem{prblm}[thm]{Problema}
\newtheorem{prcp}[thm]{Princípio}

\theoremstyle{remark}
\newtheorem*{rem}{Nota}

\geometry{
 a4paper,
 left=3cm,
 top=2cm,
 right=2cm,
 bottom=2cm
}

\begin{document}
\begin{center}
    \textbf{\large UNIVERSIDADE FEDERAL DO RIO GRANDE DO SUL} \\
    \textbf{\large INSTITUTO DE MATEMÁTICA E ESTATÍSTICA} \\
    \vspace{3cm}
    \textbf{\Large ANÁLISE REAL} \\
    \vspace{1cm}
    \textbf{\large NOTAS DE AULA E SOLUÇÕES PARA EXERCÍCIOS SELECIONADOS} \\
    \vspace{3cm}

    \textbf{LEONARDO SUAREZ}

    \vfill
    \textbf{\large Porto Alegre, RS} \\
    \textbf{\large \today}
\end{center}

\newpage
\tableofcontents
\newpage

\section{Espaços Vetoriais}
\section{Subespaços}
\section{Bases}
\section{Transformações Lineares}
\section{Produto de Transformações Lineares}
\section{Núcleo e Imagem}
\section{Soma Direta e Projeção}
\section{A Matriz de uma Transformação Linear}
\section{Eliminação}
\section{Produto Interno}
\section{A Adjunta}
\section{Subespaços Invariantes}
\section{Operadores Auto-Adjuntos}
\section{Operadores Ortogonais}
\section{Operadores Normais}
\section{Pseudo-inversa}

\end{document}